\chapter{理论引擎：广义黏滞的物理基础与实证路线}
\markboth{附录 A\quad 理论引擎}{}

本书正文以机制叙述为主，公式只在必要处出现。这个附录把散在各章的物理依据集中起来，交代它们的学术来源、数学形式、可检验的推论和适用限度。它的目的是让前面那些看起来像是类比的表述，落回到可以被检查和被反驳的位置上。

\section{学术谱系与理论增量}

先说清楚本书的想法从哪里来，以及相对于已有论述增加了什么。关于技术与社会之间速度不匹配的判断并不新，它有一批各自独立、只在局部重叠的来源，本书的贡献不在于提出这个现象，而在于把它统一到一组可以参数化的约束上。

第一条来源是比特与原子的区分。Nicholas Negroponte 在《数字化生存》里用 bit 和 atom 的对立来描述数字化的边界，指出可以被数字化的东西和必须保持物质形态的东西遵循不同的规律。这个区分被广泛引用，但在他那里主要用于描述商业形态的变化，没有给出两种介质为什么不同的机制说明。\cite{negroponte1995being} 本书把它推进了一步：比特、原子与生命的差别不是形态差别，而是吸收机制差别：信息被转译为扩散，物质被转译为做功，生命被转译为反馈。三种介质的中文名称与英文对照为比特（Bits）、原子（Atoms）、生命（Gene），最后一个取遗传信息之意，而不指细胞里那段具体的遗传物质。这个推进带来一个可检验的推论：同一项智能技术进入三种介质时，其效果的时间曲线形状不同，扩散型是单调趋近，做功型受通量上限约束，反馈型会出现非线性甚至反向。

第二条来源是文化滞后。William F. Ogburn 在 1922 年提出，社会的物质文化变化快于非物质文化，两者之间的时间差造成失调，他把这个现象称为文化滞后。\cite{ogburn1922social} 这个概念抓住了问题的形状，但缺乏度量：什么是“快”，什么是“慢”，差别多大会产生失调，都没有给出可操作的量。本书用时间刚性 $T$ 给出度量，把“滞后”还原为两个过程的弛豫时间之比，从而把定性的判断变成可以比较的量。

第三条来源是生产率悖论。Robert Solow 在 1987 年的一句话被反复引用：计算机时代无处不在，唯独在生产率统计里看不到。\cite{solow1987watch} 此后三十年的研究大致收敛到一个结论：信息技术确实提高了生产率，但存在显著的滞后，而且滞后的长短取决于组织是否完成了配套的流程重组。这个结论与本书的拓扑耦合 $G$ 是同一件事的两种说法：技术进入系统后，产出不会立刻变化，因为系统需要重新布线，而重新布线的速度由组织结构和制度安排决定，不由技术本身决定。

第四条来源是成本病。Baumol 和 Bowen 在 1966 年对表演艺术的研究是这一说法的来源：在这类活动上，产出受制于现场的人力投入，效率提升的空间有限，当其他部门的工资随生产率上升时，它们的单位成本会被动抬高；这个概念后来被推广到教育、医疗与照护等同样难以提速的领域。\cite{baumol1966performing} 它抓住的是价格信号，不匹配会以相对价格分化的形式被经济系统记录下来，而这一点可以直接从产业统计里观察到。本书接上的是它没有回答的那一环：这些位置为什么难以被其他机制替换。答案是时间刚性与拓扑耦合，它们不随算力下降，因此这类相对价格不会因为生成能力普及而回落。两种判断落在同一批对象上，依据不同，合起来比单独使用更可靠。

这些来源分别处理数字化、时间差、组织互补和相对成本，已有各自的度量与模型。这里提出的工作，是用资源、时间与耦合组织同一条现实转换链，继续把验证和责任列入比较。评价标准是它能否增加解释、区分方案并改善决定。

相对于这四条来源，本书的理论增量可以归结为三点。一是把三种阻力统一成一组参数 $(R, T, G)$，使不同类型的黏滞可以在同一个框架下比较和置换。二是给出介质的分类，说明同样是智能输入，进入不同介质后被吸收的方式不同，因此不能用同一套加速逻辑处理。三是把验证与承责纳入约束，作为与物理阻力并列的一类黏滞。第三点是本书与前人差别最大的地方，也是最需要被检验的地方。

这四条来源各自还带着一层本书没有采纳的判断，把它们分清楚可以避免误读。文化滞后理论在提出者那里常常带有对传统的肯定意味，而本书只取其中的时间差结构，不取任何关于快慢孰优的结论；成本病理论常被用来论证某些部门效率低下，而本书只取相对价格上升的机制，不取对部门本身的评价；生产率悖论在早期被当作技术无效的证据，后来被当作滞后即将结束的乐观预期，本书对两种读法都不采纳，只保留技术效果需要配套重组这一条；比特与原子的区分常被用来判断哪些行业会被数字化，本书不取这一层预测用法，只取两种介质吸收机制不同这一条。做这样的切割是有代价的：能说的话变少了，能被误用的地方也变少了，而一个框架的可信度主要取决于后一条。

需要补一句关于增量的性质的话。这三点都属于形式上的统一，不属于新机制的发现：参数化的价值在于让不同领域的观察可以在同一张表上比较，而不在于提出了某个此前无人知道的机制。因此对本书的评价标准也应当相应调整：不应当问它是否发现了新的物理，而应当问它给出的这组参数是否真的能区分不同的情形，以及用它做出的推论是否可以被证伪。
